\documentclass[11pt,a4paper]{article}
\usepackage[margin=24mm]{geometry}
\usepackage{hyperref}
\usepackage{fontspec}
\setmainfont{LibertinusSerif-Regular.otf}[BoldFont=LibertinusSerif-Bold.otf,ItalicFont=LibertinusSerif-Italic.otf]
\setsansfont{FiraSans-Regular.otf}[BoldFont=FiraSans-Bold.otf,ItalicFont=FiraSans-Italic.otf]
\setmonofont{cmuntt.otf}
\usepackage{polyglossia}
\setdefaultlanguage{russian}
\usepackage{amsmath,booktabs,pgfplots}
\pgfplotsset{compat=1.18}
\definecolor{pine}{HTML}{006C64}
\definecolor{ink}{HTML}{183848}
\definecolor{copper}{HTML}{B75529}
\newcommand{\artm}{\href{https://github.com/bigartm/bigartm}{BigARTM}}
\newcommand{\topicjev}{\href{https://github.com/anatems1/TopicJev/tree/e271339027555160e1a9dc2efa8fc6e71b1a5475}{TopicJev}}
\newcommand{\jev}{\href{https://typesafe.ai/blog/introducing-system-one-models-and-jev}{Jev}}
\newcommand{\laya}{\href{https://huggingface.co/convaiinnovations/laya-typed-decisions}{Laya}}
\newcommand{\luna}{\href{https://developers.openai.com/api/docs/guides/decisions}{OpenAI Decisions}}
\makeatletter\let\rowinput\@@input\makeatother
\newcommand{\basetable}{\begin{tabular}{llrrrr}\toprule Корпус & Метод & PPL $\downarrow$ & NPMI & TD & AMI $\pm s$\\\midrule\rowinput baseline-rows.tex \bottomrule\end{tabular}}
\newcommand{\dectable}{\begin{tabular}{llrrrr}\toprule Корпус & Модель & AMI до & AMI после & Охват & Согласие \\\midrule\rowinput decision-rows.tex \bottomrule\end{tabular}}
\newcommand{\pplplot}{\begin{tikzpicture}\begin{axis}[width=.83\linewidth,height=4.8cm,ybar,bar width=13pt,ymin=0,ymax=110,ytick={0,25,50,75,100},ylabel={\% от PPL униграммы},xtick={1,2,3},xticklabels={PLSA,Сглаживание,Сглаж. + декор.},xmin=.5,xmax=3.5,axis lines*=left,legend columns=2,legend style={at={(.5,1.02)},anchor=south,draw=none},ymajorgrids,grid style={gray!22},tick label style={font=\small}]\addplot[fill=pine,draw=pine] table[x=method,y=ag]{ppl.dat};\addplot[fill=copper,draw=copper] table[x=method,y=ng]{ppl.dat};\legend{AG News,20 Newsgroups}\end{axis}\end{tikzpicture}}

\hypersetup{colorlinks=true,linkcolor=pine,urlcolor=pine,pdftitle={Проверка тематических назначений}}
\setlength{\parindent}{0pt}\setlength{\parskip}{5pt}
\newcommand{\head}[1]{\par\smallskip\textbf{#1}\par}
\begin{document}
{\sffamily\color{pine}МФТИ · Вероятностное тематическое моделирование}\par
{\LARGE\sffamily\bfseries Проверка тематических назначений\\ с помощью моделей принятия решений}\par
Фанис Хафизов \hfill 26 ноября 2026
\section{Что означает «документ относится к теме»}
В модели мешка слов документ является смесью тем. Проверка словесного описания ставит другой вопрос: подходит ли данному тексту одна из предложенных интерпретаций? В новости о росте прибыли производителя чипов могут одновременно встречаться финансовая и технологическая лексика. Описания «результаты компаний» и «разработка новых процессоров» уточняют, какой факт мы хотим классифицировать. Исследуем, приближает ли такой выбор тематические назначения к известным рубрикам двух корпусов.

\head{Постановка и целевая функция}
Для словаря $W$, документов $D$ и $K$ тем выбираем стохастические матрицы $\phi_{wt}=p(w\mid t)$ и $\theta_{td}=p(t\mid d)$. В PLSA правдоподобие объясняет частоты $n_{dw}$ через $p(w\mid d)=\sum_t\phi_{wt}\theta_{td}$. ARTM добавляет в критерий взвешенные регуляризаторы:
\[
 \mathcal L=\sum_{d,w}n_{dw}\log\sum_t\phi_{wt}\theta_{td}+\sum_j\tau_j R_j(\Phi,\Theta).
\]
E-шаг распределяет наблюдённую частоту слова между темами пропорционально $\phi_{wt}\theta_{td}$. M-шаг обновляет нормированные тематические счётчики с поправкой на регуляризаторы. Положительное сглаживание поддерживает малые компоненты. Декорреляция понижает вес слов, которым несколько тем дают одновременно высокую вероятность. Описание темы формируем из столбца $\Phi$, исходное жёсткое назначение получаем как $\arg\max_t\theta_{td}$.

\head{От Jev к открытой реализации}
TypeSafe описывает \jev{} как модель типизированных решений. Запрос содержит состояние и вопрос. Для выбора из списка модель возвращает распределение оценок и выбранный вариант. Авторы сообщают о параллельном вычислении ответов, двухэтапной обработке больших списков и обучении RLCD. Этого описания недостаточно, чтобы восстановить точную архитектуру Jev.

Открытая \laya{} даёт конкретный пример. Англоязычный чекпойнт использует ModernBERT-large, 421 млн параметров и контекст в 1024 токена. На вход подаются вопрос, варианты с маркерами и документ. Двунаправленное внимание связывает маркеры вариантов со свидетельствами в тексте. К скрытым представлениям прибавляется вектор типа вопроса. После двух дополнительных слоёв Transformer общий скорер преобразует представление каждого маркера в скаляр. Softmax с температурой даёт распределение по предъявленным вариантам. Значение зависит от вопроса и его альтернатив. Параметр $\theta_{td}$ имеет отдельную интерпретацию в модели порождения слов.
\newpage
\section{Экспериментальный протокол}
\head{Корпуса и подготовка данных}
\href{https://github.com/mhjabreel/CharCnn_Keras/tree/master/data/ag_news_csv}{AG News} содержит заголовки и описания новостей четырёх рубрик. \href{https://scikit-learn.org/stable/datasets/real_world.html#the-20-newsgroups-text-dataset}{20 Newsgroups} содержит сообщения двадцати дискуссионных групп. Использованы исходные train/test, для 20 NG — разбиение по времени. Убраны заголовки писем, цитаты и подписи средствами scikit-learn. Это снижает возможность угадывать группу по служебным полям. Из сообщений с пятью и более буквенными словами взяты стратифицированные подвыборки по 4000/800 документов. После ограничения словарём требовался хотя бы один известный токен на обучении и два на тесте. Итоговые размеры — 4000/800 и 3998/799.

Словарь строился на обучении по английским словам длиной от трёх букв. Исключены английские стоп-слова scikit-learn, термы с документной частотой меньше четырёх или выше 60\%. Максимум — 10\,000 признаков. Остались 4171 слово на AG News и 9563 на 20 NG. Рубрики используются для стратификации, задания $K$ и внешнего оценивания. Тексты рубрик в запрос верификатора не входят.

\head{Методы и воспроизводимость}
Использован официальный CLI \artm{} (коммит \texttt{47e37f98}). Для каждого корпуса построены PLSA; ARTM со сглаживанием $\tau_\Theta=0{,}1$, $\tau_\Phi=0{,}01$; та же модель с дополнительно включённым DecorrelatorPhi ($\tau=5000$). Это фиксированные конфигурации. Использованы инициализации 11, 23 и 47, один поток, 25 внешних и 10 внутренних итераций. Тестовые векторы получены при фиксированной обученной $\Phi$.

\head{Как сопоставляются метрики}
Для document completion каждое вхождение известного слова независимо попадает в наблюдаемую часть с вероятностью $1/2$. Документы с пустой половиной исключены только из этой метрики. Оцениваем
\[
 \mathrm{PPL}=\exp\bigg(-\frac1N\sum_{d,w}n_{dw}^{\rm hidden}
  \log\big[0{,}99\,p(w\mid\theta_d^{\rm obs})+0{,}01\,p_{\rm train}(w)\big]\bigg).
\]
Общая для методов униграммная примесь защищает от нулевой вероятности скрытого слова. В тесте использованы 7192 скрытых токена из 800 новостей и 25\,977 из 794 сообщений. Результат относится к зафиксированному словарю; слова вне него исключены.

NPMI@10 — среднее по темам и парам топ-слов значения $\log[p(u,v)/(p(u)p(v))]/[-\log p(u,v)]$. Событие — присутствие в одном обучающем документе. Для пары с нулевой совместной частотой берём предельное значение $-1$. TD@10 равна числу уникальных топ-слов, делённому на $10K$. AMI сравнивает разбиение по ведущим темам и внешним меткам с поправкой на случайное совпадение. Метки рубрик служат приближением к тематической близости. Эти метрики оценивают разные свойства и зависят от предобработки.
\newpage
\section{Результаты BigARTM}
\begin{center}\small\basetable\end{center}
Средние по трём инициализациям. Для AMI указан разброс между запусками (выборочное стандартное отклонение). NPMI, TD и AMI предпочтительны выше, PPL — ниже.

\begin{center}\pplplot\end{center}
Частотная модель из обучающего корпуса имеет PPL 2172 на AG News и 4076 на 20 NG. Тематические модели превосходят её во всех 18 запусках. Максимальный модуль контрольной AMI при перестановке назначений меньше 0,005. Во всех запусках на тесте представлены все $K$ ведущих тем. В коде также проверены неотрицательность и нормировка $\Phi$ и тестовых $\Theta$.

Сглаживание уменьшает среднюю PPL примерно на 10\% относительно PLSA на AG News и на 4\% на 20 NG. Изменения AMI имеют порядок разброса между инициализациями. Для выбранного веса декорреляции средние результаты близки к одному сглаживанию. Числа нельзя напрямую сравнивать с публикациями, использующими полный корпус, иной словарь, оконную когерентность или другой способ разбиения токенов.

\head{Содержательная проверка тем}
У сглаженной модели 20 NG (инициализация 11) тема 3 содержит \textit{space, launch, satellite, nasa, technology, program, lunar, orbit, shuttle, data}. На AG News тема 2 объединяет \textit{game, season, iraq, people, team, win, night, killed, said, sunday}. Второй пример показывает, почему проверка по перечню слов может быть трудной. Даже при положительной AMI встречаются смешанные темы и признаки источника сообщения. В одном сравнении зафиксирована одна инициализация, поэтому результаты проверки нельзя обобщить на все три.
\newpage
\section{Проверка тем через TopicJev}
В примере авторов \topicjev{} проверяет назначения BERTopic по словесным описаниям тем. Мы передали десять топ-слов BigARTM для каждой темы. Идентификатор варианта содержит два первых слова и индекс темы. Сглаживание и seed 11 зафиксированы для обоих корпусов. Проверяем одни и те же 100 тестовых документов корпуса, поровну из каждой рубрики. Вход ограничен первыми 2400 символами; дополнительно действует токенное окно Laya. Ассоциации «рубрика — тема» в модель не передаются.

Штатная реализация TopicJev добавляет к $K$ содержательным вариантам Other. Если победил содержательный вариант, его вероятность должна строго превышать $2/K$. Иначе результат также Other. Доля принятых документов — охват. В AMI на всех документах Other рассматривается как отдельный кластер. Согласие — доля совпадений с ведущей темой исходной модели, включая отказы в знаменатель.
\begin{center}\small\dectable\end{center}

\head{Выбор Other и вклад порога}
\begin{center}\small
\begin{tabular}{llrrrrr}\toprule
Корпус & Модель & Other$^1$ & Порог$^2$ & AMI$^3$ & До$^4$ & После$^4$\\\midrule
\rowinput threshold-rows.tex \bottomrule\end{tabular}\end{center}
{\small $^1$ Победил класс Other. $^2$ Дополнительные отказы из-за порога. $^3$ AMI всех документов при одном argmax, включая Other. $^4$ Парная AMI только на принятых документах.}

На AG News фильтр оставил 23 новости и две содержательные темы из четырёх. AMI выросла на этом отобранном множестве, но снизилась на всех 100 документах. Для новости AG News с ID 7370 о переговорах Fiat и GM ведущая тема BigARTM содержит \textit{company, year, million}. Четыре содержательных варианта Laya получили приблизительно 0,239, 0,234, 0,045 и 0,241; остаток достался Other. Даже победитель набрал меньше порога 0,5.

На 20 NG все десять отказов возникли при прямом выборе Other. Порог $2/20=0{,}1$ дополнительных документов не отсеял. При загрузке чекпойнта Laya библиотека ограничила температуру корзины запросов с 11 и более вариантами снизу значением 0,5. Поэтому численные оценки для 20 NG особенно важно проверять на отдельной калибровочной выборке.

\newpage
\section{Как меняется охват при отказе}
\begin{center}
\begin{tikzpicture}
\begin{axis}[xbar stacked,width=.9\linewidth,height=4cm,bar width=13pt,symbolic y coords={ng,ag},ytick={ng,ag},yticklabels={20 Newsgroups,AG News},xmin=0,xmax=100,enlarge y limits=.7,axis lines*=left,xlabel={Документов из 100},legend columns=3,legend style={at={(.5,1.08)},anchor=south,draw=none,font=\small}]
\addplot[fill=pine,draw=pine] table[x=accepted,y=dataset]{routing.dat};
\addplot[fill=copper,draw=copper] table[x=direct,y=dataset]{routing.dat};
\addplot[fill=gray!45,draw=gray!45] table[x=threshold,y=dataset]{routing.dat};
\legend{Тема принята,Выбран Other,Отсечено порогом}
\end{axis}
\end{tikzpicture}\end{center}
Распределение решений Laya показывает два механизма отказа. Для четырёх рубрик AG News порог 0,5 переводит в Other ещё 43 документа. На 20 NG переход от одного argmax к штатному порогу оставляет те же назначения.

\head{Документы, для которых назначения совпали}
В новости AG News с ID 1288 о выборе стартового квотербека команды Chargers обе модели дали тему 2. Оценка Laya 0,629 превысила порог. В этой же теме среди топ-слов соседствуют спортивная и военная лексика. Согласие по конкретной новости показывает, что часть смешанной темы сохраняет понятное значение.

Обе модели отнесли к теме 19 два сообщения. В первом спрашивают о накопителе для Centris 650, во втором продают Macintosh с жёстким диском. Оценки Laya составили 0,963 и 0,723. Эталонные рубрики различаются: обсуждение оборудования Macintosh и частные объявления. Оба сообщения описывают оборудование; рубрика объявления учитывает также цель автора. Даже осмысленная предметная тема может пересекать границы эталонных групп.

\head{Границы вывода и следующий замер}
Ухудшение AMI характеризует соответствие внешней рубрикации при этих описаниях тем. Для оценки точности именно словесных описаний нужна независимая разметка пар «документ, тема». Для \luna{} подготовлен адаптер штатного JevEntail с моделью \texttt{gpt-6-luna}. Меняется удалённый вызов; сохраняются варианты, Other, порог и тестовые ID. Полный ответ и usage сохраняются локально. В таблицы попадают только завершённые измерения.

\small\textbf{Материал курса.} К. В. Воронцов, лекции 1, 2, 3, 5 и 7 весны 2026: постановка, совместная встречаемость слов, энкодеры, регуляризация и качество. Для document completion см. лекцию 7, слайд 8. Для архитектуры см. \href{https://laya.studio/learn/typed-decisions}{описание маркеров Laya}; для порога — \href{https://github.com/anatems1/TopicJev/blob/e271339027555160e1a9dc2efa8fc6e71b1a5475/topicjev/entail/base.py}{реализацию EntailResults}.
\end{document}
